\section{MoE systems side：expert parallelism 与通信}

前面的 routing 公式只说明“某个 token 应该去哪里”，还没有说明 token 真正如何跨 GPU 到达 expert。本节把 MoE 还原成运行时流水线：router 在本地计算目的地，dispatch 按 expert 重排 token，all-to-all 把 activation 发送到持有对应权重的设备，expert GEMM 完成计算，随后 combine communication 把结果送回原序列位置。每一步都可能引入 padding、同步和尾延迟，因此算法层的稀疏 FLOPs 只有在系统流水线也高效时才会变成 wall-clock 收益。



\begin{figure}[H]
\centering
\includegraphics[width=0.88\textwidth]{slide-043.jpg}
\caption{Slide 44：expert parallelism；不同设备保存不同 experts。}
\figsource{CS336 Spring 2026 Lecture 4 官方幻灯片。}
\end{figure}

\noindent\textbf{展开说明：}MoEs parallelize nicely – Each FFN enables additional kinds of parallelism。这里的“nice”有前提：expert weights 可以天然分片，但 token 必须按路由结果发送到远端设备，再把输出发回；计算分片简单，通信调度并不简单。吞吐取决于每卡 token 数是否均衡、all-to-all 是否与 expert GEMM 重叠，以及小 expert matrix 是否仍能高效利用 GPU。

\begin{knowledgebox}{讲义补充：源 Slide 44 应该怎么看}
这页是机制或证据页。先看它比较的是 compute、参数量、routing、负载均衡还是通信；再看结论支持的是质量提升、训练速度、推理成本还是系统复杂度。MoE 和 attention alternatives 的图常把模型结构和系统代价混在一起，读图时要分清“算法机制”和“硬件执行成本”。
\end{knowledgebox}


\begin{figure}[H]
\centering
\includegraphics[width=0.88\textwidth]{slide-044.jpg}
\caption{Slide 45: Training MoEs – the systems side}
\figsource{CS336 Spring 2026 Lecture 4 官方幻灯片。}
\end{figure}

\noindent\textbf{展开说明：}MoE routing allows for parallelism, but also some complexities

\begin{knowledgebox}{读图：Slide 45 应该怎么看}
这页是机制或证据页。先看它比较的是 compute、参数量、routing、负载均衡还是通信；再看结论支持的是质量提升、训练速度、推理成本还是系统复杂度。MoE 和 attention alternatives 的图常把模型结构和系统代价混在一起，读图时要分清“算法机制”和“硬件执行成本”。
\end{knowledgebox}


\begin{figure}[H]
\centering
\includegraphics[width=0.88\textwidth]{slide-045.jpg}
\caption{Slide 46: MoE parallelism and architecture modifications}
\figsource{CS336 Spring 2026 Lecture 4 官方幻灯片。}
\end{figure}

\begin{figure}[H]
\centering
\includegraphics[width=0.88\textwidth]{slide-046.jpg}
\caption{Slide 47：MoE stochasticity；路由会给训练和复现引入额外随机性。}
\figsource{CS336 Spring 2026 Lecture 4 官方幻灯片。}
\end{figure}

\begin{warningbox}{Expert parallelism 不是免费扩容}
MoE 把 dense FFN 的本地 GEMM 变成 dispatch、all-to-all、expert compute、combine 四阶段。若 batch 太小、路由过于偏斜或 experts 切得太细，通信和 launch overhead 会主导；若为了平衡加入 padding，又会重新花掉被稀疏性省下的 FLOPs。系统验收应报告每阶段时间、发送字节、capacity overflow 和专家利用率。
\end{warningbox}

\noindent\textbf{展开说明：}New ideas from Nemotron 3 – down-projecting the activations to reduce communication

\begin{knowledgebox}{读图：Slide 46 应该怎么看}
这页是机制或证据页。先看它比较的是 compute、参数量、routing、负载均衡还是通信；再看结论支持的是质量提升、训练速度、推理成本还是系统复杂度。MoE 和 attention alternatives 的图常把模型结构和系统代价混在一起，读图时要分清“算法机制”和“硬件执行成本”。
\end{knowledgebox}



\noindent\textbf{展开说明：}MoEs can have additional stochasticity beyond normal models..

\subsection{本章小结}
本节的共同线索是 selective computation：通过结构选择让 token 不必总是访问所有历史或所有参数。但选择性越强，越需要额外机制保证信息流、负载均衡和训练稳定。

